\section{Earlier cases and limits of the comparison}\label{sec:recovery}

\subsection{Extension families}
For the bipartite two-step families of
\cite[Theorem~1.1]{wuebben1}, take $m_i=1$ and restrict to the
specified extension subspace. There are no composable products,
so $a=x$ and $\kappa=0$ on that subspace. The commutator form
\eqref{eq:finite-spectrum} becomes the weighted graph form.
If $\sigma$ is the total normalized balanced edge energy, then
\[
 \|x\|_t^2=\Vol(\omega_t)\sigma .
\]
Proposition~\ref{prop:moving} gives the earlier $O(\sigma)$
metric and spectral rates without assuming $|t|=O(\sigma)$.
The statement is first made at the small balanced representative
and then pulled back by the original harmonic identification and
balancing gauge. Metric-order inequalities survive this pullback
even when that original gauge diverges. The fixed extension,
prescribed nonzero components and its proven stability cone
remain the hypotheses for the earlier stability criterion.
Tree determinant identities require their separate finite
graph argument.

Allowing arbitrary fixed multiplicities gives the complete
complex algebra of \cite[Theorem~1.1 and Section~5]{wuebben4},
with dimension $\sum_i m_i^2-1$. The mass and energy agree after
a common volume normalization. Simplicity and the fixed stable
multiplicity orbit are still needed for an entirely positive
trace-free small spectrum and a unique normalized metric.
The present theorem also includes nonsimple polystable points;
it does not give an HE metric to a nonsplit semistable bundle.

For the adjacent-arrow three-step family of
\cite[Section~3]{wuebben5}, the only quadratic product lies in
the endpoint block. Every product involving that block vanishes.
The complete canonical equation therefore terminates:
\[
 a=x-d_t^*G_{2,t}(x\wedge x).
\]
Its integrability equations are precisely the displayed cup-product
relations of that family. This is a consequence of its triangular
algebra, not a truncation applicable to arbitrary quivers.
The normalized harmonic energy is the sum of the balanced edge
weights, so Proposition~\ref{prop:moving} preserves the earlier
quadratic relative rate, including arbitrary positive ratios.
The distinct line factors, adjacent-arrow subspace and stable
connected supports remain part of the graph interpretation.

\subsection{Full local moduli and K3 coordinates}
The complementary-minimization result
\cite[Section~8]{wuebben6} is an imported principle in a more
general compact-complex formulation. The harmonic arguments here
verify its hypotheses in the moving presentation. The
$83$-coordinate wall chart of that paper instead uses a fixed
reference Hodge contraction. Its product geometry gives the
required orthogonality in the current norms: each harmonic Hom
block has one differential leg, and changing the polarization
multiplies the pairing on that block by a constant.
Thus the reference-orthogonal correction remains orthogonal.
Its quadratic quotient estimate is not obtained by treating an
arbitrary change of analytic chart as an isometry.
At a fixed endpoint the estimate has its endpoint-dependent
constant $C_l$; no uniformity as $l\to0$ follows from the
fixed-background theorem. The additional endpoint and limiting-metric
conclusions retain their original hypotheses and proofs.

The full model of \cite[Sections~2, 4--6]{wuebben7} is
\[
 E_0=F_1\oplus(F_2\otimes\C^2)\oplus F_3\oplus F_4
                    \quad\hbox{on }(\PP^1)^3 .
\]
Its complete harmonic recursion has $119$ coordinates, and its
parallel trace-free algebra has dimension
$1^2+2^2+1^2+1^2-1=6$. The line-factor product metrics are HE for
the given product polarizations. The normalized energy $q_t$
equals $\|x\|_t^2/\Vol(\omega_t)$, so the same $O(q_t)$
stationary metric, spectral and quotient rates follow.
All six modes, including zeros, are retained.
The general finite correction covers nonsimple flat-balanced
points without a scalar-centralizer assumption.
The explicit weak coefficient, geometric curve and
parameter-dependent length assertions remain additional results
of that example.

For the K3 comparisons in
\cite[Sections~4, 6]{wuebben13} and
\cite[Section~3]{wuebben14}, one must distinguish canonical
harmonic coordinates from the independently established
convergent equivariant formality coordinates. Suppose $F$ maps
the quadratic locus to the full Kuranishi locus and
$F(x)=x+O(\|x\|^2)$. Then $DF=I+O(r)$, and exact equivariance gives
\begin{equation}\label{eq:formality-equivariance}
 DF_xC_x\xi=C_{F(x)}\xi.
\end{equation}
Applying the derivative estimate to $C_x\xi$ itself yields a
relative energy comparison even in a very weak orbit direction.
The pullback ambient form has error $O(|t|+r)$, giving that
rate in the straightened coordinates. If the verified symmetry
makes $F$ odd, then $DF=I+O(r^2)$, recovering $O(|t|+r^2)$.
For the real equation one must also apply finite balancing to
the pulled-back moment form. A nonlinear analytic coordinate
isomorphism need not preserve a flat moment level.
For quotient norms, \eqref{eq:formality-equivariance} identifies
the orbit directions, while correction derivatives disappear
under minimization. The formality hypotheses, homogeneous
rescaled limits and intrinsic distance results of those papers
are retained, not inferred from a general obstructed germ.

The rank-three K3 example uses the opposite sign at its
multiplicity-two vertex. In its notation
$\mu_{13}=A^*A-BB^*$ and $\tau(t)=-28\pi t/3$.
Equation~\eqref{eq:moment} instead gives
\[
 \mu_1=-\mu_{13},\qquad \zeta_1(t)=28\pi t/3=-\tau(t).
\]
The second scalar equation follows by trace. This identifies the
same real level set with the curvature convention $+iF^0$.

For all these comparisons, common division of zero-form mass and
one-form energy by $\Vol(\omega_t)$ leaves eigenvalues unchanged.
The normalized graph convention uses edge weights
$\|x_a\|_t^2/\Vol(\omega_t)$ and vertex masses $r_i$.
The unnormalized convention uses $\|x_a\|_t^2$ and
$\Vol(\omega_t)r_i$. No dimension-free statement is made.

\subsection{Boundaries}
The relative estimates require more than small absolute errors.
For instance the nonnegative matrices
\[
 \diag(\varepsilon^2,r^2),\qquad
 \diag(\varepsilon^2+r^4,r^2)
\]
differ in norm by $r^4$, while the ratio of their least eigenvalues
is $1+r^4/\varepsilon^2$. It is unbounded when
$\varepsilon/r^2\to0$. The moment projection estimate
\eqref{eq:projection-error} and the spectral energy bound
\eqref{eq:relative-block} prevent this failure in their respective
arguments. Both use exact geometric identities.

The theorem is local near a fixed smooth locally free HE
background with positive complementary elliptic gaps.
It does not supply harmonic spaces or such gaps for a singular
graded sheaf, or uniformly through collapse of the background
geometry. The singular-flow obstruction
\cite[Theorem~1.1 and Corollary~1.2]{wuebben10}
has a precise different scope. It concerns a locally free rank-two
extension on a compact K\"ahler surface,
\[
 0\longrightarrow L\longrightarrow V\longrightarrow I_WM
                    \longrightarrow0,\qquad
 \deg L=\deg M,
\]
where $W$ is a nonempty zero-dimensional subscheme.
Its induced line metric coefficient $u(T)$ along HYM flow
satisfies $\sqrt T\,u(T)\to0$ and
$\log u(T)+\tfrac12\log T\to-\infty$ on every nonempty compact
subset away from $\operatorname{Supp}W$.
This excludes the stated bounded metric-order and ordinary
matrix-log comparisons with the categorical $T^{\pm1/2}$ model
for fixed compatible splittings. The stationary bubbling result
in that work is separately restricted to its toric example.
Neither assertion is a general impossibility theorem for
stationary local reductions at singular sheaves.

Finally, stationary HE reconstruction and zero-form spectral
estimates alone give no HYM trajectory, matching of a prescribed
initial metric, asymptotic phase, differentiated response estimate,
global metric completion or homogeneous blow-up limit.
Those questions require additional dynamical or geometric
arguments. The full analytic obstruction, genuine reduced
tangent spaces and fixed-background dependence in
Theorem~\ref{thm:main} specify the local stationary result.

